\scanpage{39}
% Pencil archive foliation: 38.
\uline{gradué} \inserted{\uline{à degrés} $\geq0$}, et $\xi$ est de degré $1$. De cette
façon, $\mathrm{M}(X)$ devient gradué [à degrés
de tous signes, $\xi$ étant une \unsure{droite}
\uline{inversible} de degré $1$]. Noter que
\[
\mathrm{M}_0(X)\simeq\mathbb{Z}^{\Sigma_0(X)}
\]
avec
\[
\Sigma_0(X)=\underset{\substack{\longrightarrow\\
\text{par }x\unsure{\rightsquigarrow|\mapsto}x.\xi}}{\lim}\Sigma_i^+(X)
\]
et
\[
\mathrm{M}^{\mathrm{pair}}(X)=
\underbrace{\mathrm{M}_0(X)}_{\text{de degré }0}
[\underbrace{\xi}_{\text{de degré }1}]
\]
\inserted{\[
\mathrm{M}^{\mathrm{impair}}(X)=
\mathrm{M}_1(X)\otimes_{\mathrm{M}_0(X)}\mathrm{M}_0(X)[\xi]
\]}
donc la structure d'anneau gradué
de $\mathrm{M}(X)$ est déterminée en termes de
la structure d'anneau de $\mathrm{M}_0(X)$, \deleted{et}
du module $\mathrm{M}_1(X)$ sur $\mathrm{M}_0(X)$, et l'application de mult.
\inserted{définie par \deleted{\unsure{\xi}} $u,v\unsure{\rightsquigarrow|\mapsto}\xi^{-1}uv$.}
\[
\mathrm{M}_1(X)\otimes_{\mathrm{M}_0(X)}\mathrm{M}_1(X)
\longrightarrow\mathrm{M}_0(X)
\]

D'autre part les objets de $D_i^+(\mathcal{M}^+(X))$
sont \deleted{ceux qui}\inserted{ceux} dont la classe dans
\deleted{$\mathcal{M}^+(X)$ sont} $\mathrm{M}^+(X)$ sont dans
$\mathbb{Z}^{D_i(\Sigma^+(X))}$.
\[
\left.
\begin{aligned}
D_i(\Sigma^+(X))&=
\coprod_{j\geq0}\deleted{\unsure{\Sigma}}\xi^j\Sigma_{(i)}^+(X)\\
&=\coprod_{\substack{j\geq0\\0\leq k\leq i}}\xi^j\Sigma_k^+(X)
\end{aligned}
\right\}\Sigma_{[i]}^+(X)
\]
